% This is part of Mes notes de mathématique
% Copyright (c) 2011-2019
%   Laurent Claessens
% See the file fdl-1.3.txt for copying conditions.

\label{CHAPGroupes}

Pour rappel, la notion de groupe est définie en \ref{DEFooBMUZooLAfbeM}.

%+++++++++++++++++++++++++++++++++++++++++++++++++++++++++++++++++++++++++++++++++++++++++++++++++++++++++++++++++++++++++++ 
\section{Somme à valeurs dans un groupe abélien}
%+++++++++++++++++++++++++++++++++++++++++++++++++++++++++++++++++++++++++++++++++++++++++++++++++++++++++++++++++++++++++++

%--------------------------------------------------------------------------------------------------------------------------- 
\subsection{Somme à valeurs dans un anneau}
%---------------------------------------------------------------------------------------------------------------------------

\begin{definition}      \label{DEFooNEVNooJlmJOC}
    Si \( f\colon \{ 0,\ldots, N \}\to A\) est une application vers un anneau \( A\), alors nous définissons la notation \( \sum_{i=0}^Nf(i)\) par récurrence de la façon suivante :
    \begin{enumerate}
        \item
            \( \sum_{i=0}^0f(i)=f(0)\),
        \item
            \( \sum_{i=0}^{k}f(i)=\sum_{i=0}^{k-1}f(i)+f(k)\).
    \end{enumerate}
\end{definition}

%--------------------------------------------------------------------------------------------------------------------------- 
\subsection{Somme à valeurs dans un groupe abélien}
%---------------------------------------------------------------------------------------------------------------------------

Si \( S\) est un ensemble fini, nous savons de la proposition \ref{PROPooJLGKooDCcnWi} qu'il existe un unique \( N\in \eN\) pour lequel il existe une bijection \( \varphi\colon \{ 0,\ldots, N \}\to S\). Cette bijection n'est à priori pas unique.

\begin{lemmaDef}[\cite{MonCerveau}]       \label{DEFooLNEXooYMQjRo}
    Soient un groupe abélien \( (G,+)\) ainsi qu'un ensemble fini \( I\) contenant \( n\) éléments. Soit une application \( f\colon I\to G \). Si \( \sigma_1,\sigma_2\colon \{1,\ldots, n \}\to I\) sont deux bijections, alors\footnote{Pour rappel, le symbole \( \sum_{i=1}^n\) est défini par \ref{DEFooNEVNooJlmJOC}.}
    \begin{equation}
        \sum_{i=1}^nf\big( \sigma_1(i) \big)=\sum_{i=1}^nf\big( \sigma_2(i) \big).
    \end{equation}
    La valeur commune est notée
    \begin{equation}
        \sum_{i\in I}f(i)
    \end{equation}
\end{lemmaDef}

La définition \ref{DEFooLNEXooYMQjRo} donne lieu à un certain nombre de remarques.
\begin{enumerate}
    \item
        Elle donne la somme sur un ensemble fini. Un problème avec les ensembles infinis (outre la convergence) est l'ordre de sommation. Si vous voulez sommer sur \( \eZ\), dans quel ordre le faire ?
    \item
        Pour aller plus loin, et sommer sur des ensembles infinis, il faut regarder la définition \ref{DefHYgkkA}. 
\end{enumerate}

\begin{proposition}     \label{PROPooJBQVooNqWErk}
    Soient un groupe abélien \( (G,+)\), un ensemble fini \( I\) est un ensemble fini, une application \( f\colon I\to G\) et une bijection \( \sigma\colon I\to I\). Alors
    \begin{equation}
        \sum_{i\in I}f(i)=\sum_{i\in I}f\big( \sigma(i) \big).
    \end{equation}
\end{proposition}

Si nous avons une application \( L\colon S\to S\), nous notons
\begin{equation}
    \sum_{s\in S}f\big( L(s) \big)=\sum_{s\in S}(f\circ L)(s).
\end{equation}
Cette façon d'écrire donne une interprétation pour la notation \( \sum_{g\in G}f(hg)\) qui arrive dans la proposition \ref{PROPooWJQQooFINSEc}. Il s'agit de considérer l'application \( L_h\) du lemme \ref{LEMooBIBFooBHxFYC}, de considérer\footnote{Le fait que \( L_h\) soit une bijection n'a pas d'importance ici.}
\begin{equation}        \label{EQooQQBEooFDOBVG}
    \sum_{g\in G}f(hg)=\sum_{g\in G}(f\circ L_h)(g)
\end{equation}
et de faire tourner la définition \ref{DEFooLNEXooYMQjRo}. La même chose tient pour définir \( \sum_{g\in G}(gh)\) à l'aide de \( R_h\).

\begin{proposition}[\cite{MonCerveau}]     \label{PROPooWJQQooFINSEc}
    Soient un groupe fini \( G\) et une fonction \( f\colon G\to A\) où \( A\) est un anneau commutatif. Alors
    \begin{equation}
        \sum_{g\in G}f(g)=\sum_{g\in G}f(gh)=\sum_{g\in G}f(hg)
    \end{equation}
    pour tout \( h\in G\).
\end{proposition}

\begin{proof}
    Nous avons une bijection \( \varphi\colon \{ 0,\ldots,  N \}\to G\) garantie par la proposition \ref{PROPooJLGKooDCcnWi}. La définition est que
    \begin{equation}
        \sum_{g\in G}f(g)=\sum_{i=0}^Nf\big( \varphi(i) \big).
    \end{equation}
    Par ailleurs, le lemme \ref{LEMooBIBFooBHxFYC} donne une bijection \( L_h\colon G\to G\) et permet de considérer la composée
    \begin{equation}
        \begin{aligned}
            \varphi'\colon \{ 0,\ldots,  N \}&\to G \\
            \varphi'=L_h\circ \varphi.
        \end{aligned}
    \end{equation}
    La proposition \ref{DEFooLNEXooYMQjRo} nous permet d'utiliser la bijection \( \varphi'\) au lieu de \( \varphi\) pour exprimer la somme \( \sum_{g\in G}\). Ensuite un jeu de notation utilisant \eqref{EQooQQBEooFDOBVG} donne
    \begin{equation}
        \begin{aligned}[]
            &\sum_{g\in G}f(g)=\sum_{i=0}^Nf\big( \varphi(i) \big)=\sum_{i=0}^Nf\big( \varphi'(i) \big)=\sum_{i=0}^N(f\circ L_h\circ \varphi)(i)\\
            &\quad=\sum_{i=0}^N(f\circ L_h)\big( \varphi(i) \big)=\sum_{g\in G}(f\circ L_h)(g)=\sum_{g\in G}f(hg).
        \end{aligned}
    \end{equation}
    En ce qui concerne \( \sum_{g\in G}f(gh)\), c'est la même chose, en utilisant \( R_h\) au lieu de \( L_h\).
\end{proof}

Tout cela nous permet de faire une somme sympathique et bien connue.
\begin{lemma}
    Soit \( n\in \eN\). Nous avons
    \begin{equation}
        \sum_{k=0}^nk=\frac{ n(n+1) }{ 2 }.
    \end{equation}
\end{lemma}

\begin{proof}
    La preuve est pratiquement immédiate par récurrence. Nous allons donner une preuve plus «constructive», qui formalise l'idée classique d'écrire la somme à l'endroit et à l'envers.


    Nous notons \( S\) la somme \( \sum_{k=0}^nk\). Le lemme \ref{DEFooLNEXooYMQjRo} dit que si \( \sigma_i\colon \{ 0,\ldots, n \}\to \{ 0,\ldots, n \}\) sont deux bijections, alors \( \sum_{k=0}^nf\big( \sigma_1(k) \big)=\sum_{k=0}^nf\big( \sigma_2(k) \big)\). Nous sommes intéressé au cas \( f(i)=i\).

    En prenant \( \sigma_1(k)=k\) et \( \sigma_2(k)=n-k\), nous avons
    \begin{equation}
        S=\sum_{k=0}^nk=\sum_{k=0}^n(n-k).
    \end{equation}
    Donc
    \begin{equation}
        2S=\sum_{k=0}^n\big( k+(n-k) \big)=\sum_{k=0}^nn=n\sum_{k=0}^n1=n(n+1).
    \end{equation}
    En divisant par deux, nous obtenons le résultat annoncé.
\end{proof}

Dans le même ordre d'idée, pour le produit.
\begin{proposition}     \label{PROPooQMUDooQQVRIe}
    Si \( E\) est un ensemble fini et si \( G\) est un groupe abélien, alors pour toute fonction \( f\colon E\to G\) et pour toute permutation \( \sigma\) de \( E\),
    \begin{equation}
        \prod_{i\in E}f(i)=\prod_{i\in E}f\big( \sigma(i) \big)
    \end{equation}
\end{proposition}

%+++++++++++++++++++++++++++++++++++++++++++++++++++++++++++++++++++++++++++++++++++++++++++++++++++++++++++++++++++++++++++
\section{Du vocabulaire dans les groupes}
%+++++++++++++++++++++++++++++++++++++++++++++++++++++++++++++++++++++++++++++++++++++++++++++++++++++++++++++++++++++++++++

%---------------------------------------------------------
\subsection{Centre, centralisateur, normalisateur}
%---------------------------------------------------------

\begin{definition}[\cite{Kropholler}]\label{defGroupeCentre}
Soit \( G\) un groupe, et \( H\) un sous-ensemble de \( G \).
\begin{enumerate}
\item
Le \defe{centre}{centre!d'un groupe} d'un groupe \( G\) est l'ensemble des éléments de \( G\) qui commutent avec tous les autres:
\begin{equation}
    Z_G=\{ g\in G\tq gh=hg\;\forall g\in G \}.
\end{equation}

\item
Si \( h\in G\) nous notons \( Z_G(h)\) le \defe{centralisateur}{centralisateur} de \( h\) dans \( G\) :
\begin{equation}
    Z_G(h)=\{g\in G\tq gh = hg\}.
\end{equation}
C'est l'ensemble des éléments de \( G\) qui commutent avec \( g\).

\item    Le \defe{centralisateur}{centralisateur} de \( H\subset G\) est
    \begin{equation}
        Z_G(H)=\{ g\in G\tq gh=hg\,\forall h\in H\};
    \end{equation}
    il contient donc tous les éléments de \( G\) qui commutent avec ceux de \( H\).
\end{enumerate}
\end{definition}

\begin{definition}\label{DEFGroupeNormalisateur}
        Si \( H\) est un sous-groupe, son \defe{normalisateur}{normalisateur} est
    \begin{equation}
        N_G(H)=\{ g\in G\tq gH=Hg \}.
    \end{equation}
\end{definition}

\begin{remark}
Les définitions précedentes ne sont utiles que si le groupe \(G \) n'est pas abélien~\footnote{Définition~\ref{DEFGroupeAbelien}.}. Si \( G \) est abélien, alors clairement~\footnote{Vérifiez-le!}, \(Z_G = G,\ Z_G(g) = G\) pour tout \(g \in G,\ Z_G(H) = G,\ H_G(H) = G\), et tous les sous-groupes sont normaux.
\end{remark}

\begin{remark}
Les définitions précédentes ont un certain lien entre elles; en effet:
  \begin{enumerate}
    \item dans la définition~\ref{DEFGroupeCentre}, le centre est le centralisateur de \(G \), et le centralisateur d'un ensemble \( H \) est l'intersection des centralisateurs des éléments \( h \) de \( H \);
    \item un sous-groupe \( H \) est normal si et seulement si \( N_G(H) = G \).
  \end{enumerate}
\end{remark}

%---------------------------------------------------------
\subsection{Exposant}
%---------------------------------------------------------

\begin{definition}  \label{DefvtSAyb}
    L'\defe{exposant}{exposant!d'un groupe} du groupe \( G\) est le plus petit entier non nul \( n\) tel que \( g^n=e\) pour tout \( g\in G\). S'il n'existe pas, nous disons que l'exposant du groupe est infini.
\end{definition}

\begin{proposition}\label{PROPooSWHHooOzqWkw}
    Si l'ensemble des ordres de tous les éléments d'un groupe est majoré, alors l'exposant du groupe est le plus petit commun multiple des ordres des éléments du groupe

    Pour un groupe fini, l'exposant est le $\ppcm$ des ordres des éléments du groupe.
\end{proposition}

Le théorème de Burnside~\ref{ThooJLTit} nous donnera un bon paquet d'exemples de groupes d'exposant fini dans \( \GL(n,\eC)\).


%++++++++++++++++++++++++++++++++++++++++++++++++++++++++++++++++++++++++++++++++++++++++++++++ 
\section{Ordre d'un élément dans un groupe fini}
%++++++++++++++++++++++++++++++++++++++++++++++++++++++++++++++++++++++++++++++++++++++++++++++

Il y a des informations en plus dans la partie sur les groupes monogènes, \ref{SECooXIHPooWVSjhT}.

\begin{theorem}[Théorème de Cauchy\cite{ooTZHGooEPFstf}]    \label{THOooSUWKooICbzqM}
    Soit un groupe fini d'ordre \( n\). Pour tout diviseur premier \( p\) de \( n\), le groupe \( G\) possède au moins un élément d'ordre \( p\).
\end{theorem}
\index{théorème!Cauchy!groupe}

Le lemme suivant indique que sous hypothèse de commutativité, l'ordre d'un élément est une notion multiplicative.
\begin{lemma}[\cite{rqrNyg}]    \label{LemyETtdy}
    Soit \( G\) un groupe et \( a,b\in G\) tels que \( ab=ba\) d'ordres respectivement \( r\) et \( s\), deux nombres premiers entre eux. Alors l'élément \( ab\) est d'ordre \( rs\).
\end{lemma}

\begin{proof}
    Étant donné que \( (ab)^{rs}=a^{rs}b^{rs}=1\), l'ordre de \( ab\) divise \( rs\). Et vu que \( r\) et \( s\) sont premiers entre eux, l'ordre de \( ab\) s'écrit sous la forme \( r_1s_1\) avec \( r_1\divides r\) et \( s_1\divides s\). Nous avons
    \begin{equation}
        a^{r_1s_1}b^{r_1s_1}=(ab)^{r_1s_1}=1,
    \end{equation}
    que nous élevons à la puissance \( r_2\) où \( r_2\) est définit en posant \(r=r_1r_2\) :
    \begin{equation}
        a^{rs_1}b^{rs_1}=1.
    \end{equation}
    Et comme \( a^{rs_1}=1\), nous concluons que \( b^{rs_1}=1\). Donc \( s\divides rs_1\). Par le lemme de Gauss \ref{LemPRuUrsD}, nous avons en fait \( s\divides s_1\). Vu qu'on a aussi \( s_1\divides s\), nous avons \( s=s_1\).

    Le même type d'argument donne \( r=r_1\), et finalement l'ordre de \( ab\) est \( r_1s_1=rs\).
\end{proof}

\begin{lemma}[\cite{Combes}]    \label{LemSkIOOG}
    Un sous-groupe d'indice \( 2\) est un sous-groupe normal.
\end{lemma}

\begin{proof}
    Si $H$ est un tel sous-groupe d'un groupe $G$, alors $G$ possède exactement deux classes à gauche par rapport à \( H\) (théorème de Lagrange~\ref{ThoLagrange}) et se partitionne donc en deux parties : \( G=H\cup xH\) avec \( x \notin H \). De même pour les classes à droite : \( G=H\cup Hx\). Puisque la classe à droite \( Hx \) n'est pas $H$, on a \( xH = Hx \), et $H$ est normal dans $G$ par la proposition~\ref{propGroupeNormal}.
\end{proof}

\begin{lemma}[\cite{NielsBMorph}]\label{PropubeiGX}
    Soit \( H\), un sous-groupe normal d'indice \( m\) de \( G\). Alors pour tout \( a\in G\) nous avons \( a^m\in H\).
\end{lemma}

\begin{proof}
    Par définition de l'indice, le groupe \( G/H\) est d'ordre \( m\). Donc si \( [a]\in G/H\), nous avons \( [a]^m=[e]\), ce qui signifie \( [a^m]=[e]\), ou encore \( a^m\in H\).
\end{proof}

\begin{proposition}[\cite{NielsBMorph}]     \label{PROPooVWVIooQzuAlA}
    Soit un groupe fini \( G\) et \( H\), un sous-groupe normal d'ordre \( n\) et d'indice \( m\) avec \( m\) et \( n\) premiers entre eux. Alors \( H\) est l'unique sous-groupe de \( G\) à être d'ordre \( n\).
\end{proposition}

\begin{proof}
    Soit \( H'\) un sous-groupe d'ordre \( n\). Si \( h\in H'\) alors \( h^n=1\) et \( h^m\in H\) par le lemme \ref{PropubeiGX}. Étant donné que \( m\) et \( n\) sont premiers entre eux, par le théorème de Bézout~\ref{ThoBuNjam}, il existe \( a,b\in \eZ\) tels que
    \begin{equation}
        am+bn=1.
    \end{equation}
    Du coup \( h=h^1=(h^m)^a(h^n)^b\). En tenant compte du fait que \( h^n=1\) et \( h^m\in H\), nous avons \( h\in H\). Ce que nous venons de prouver est que \( H'\subset H\) et donc que \( H=H'\) parce que \( | H' |=| H |=| G |/m\).
\end{proof}

\begin{normaltext}
    Notons que cette proposition ne dit pas qu'il existe un sous-groupe d'ordre \( n\) et d'indice \( m\). Il dit juste que s'il y en a un et s'il est normal, alors il n'y en a pas d'autres.
\end{normaltext}

\begin{lemma}       \label{LemqAUBYn}
    L'ensemble des ordres\footnote{Définition \ref{DEFooKWBCooMlmpCP}.} d'un groupe commutatif est stable par PPCM\footnote{Définition \ref{DefrYwbct}.}.

    Autrement dit, si \( x\in G\) est d'ordre \( r\) et si \( y\in G\) est d'ordre \( s\), alors il existe un élément d'ordre \( \ppcm(r,s)\).
\end{lemma}

\begin{proof}
    Soit \( m=\ppcm(r,s)\). Afin d'écrire \( m\) sous une forme pratique, nous considérons les décompositions en facteurs premiers de \( r\) et \( s\) :
    \begin{subequations}
        \begin{align}
            r&=\prod_{i=1}^kp_i^{\alpha_i}\\
            s&=\prod_{i=1}^kp_i^{\beta_i}
        \end{align}
    \end{subequations}
    où \( \{ p_i \}_{i=1\ldots k}\) est l'ensemble des nombres premiers arrivant dans les décompositions de \( r\) et de \( s\). Si nous posons
    \begin{subequations}
        \begin{align}
            r'&=\prod_{\substack{i=1\\\alpha_1>\beta_i}}^kp_i^{\alpha_i}\\
            s'&=\prod_{\substack{i=1\\\alpha_i\leq \beta_i}}^kp_i^{\beta_i},
        \end{align}
    \end{subequations}
    alors \( \ppcm(r,s)=r's'\) et \( r'\) et \( s'\) sont premiers entre eux. L'élément \( x^{r/r'}\) est d'ordre \( r'\) et l'élément \( y^{s/s'}\) est d'ordre \( s'\). Maintenant nous utilisons le fait que \( G\) soit commutatif et le lemme~\ref{LemyETtdy} pour conclure que l'ordre de \( x^{r/r'}y^{s/s'}\) est \( r's'=m\).
\end{proof}

%+++++++++++++++++++++++++++++++++++++++++++++++++++++++++++++++++++++++++++++++++++++++++++++++++++++++++++++++++++++++++++
\section{Groupe dérivé}
%+++++++++++++++++++++++++++++++++++++++++++++++++++++++++++++++++++++++++++++++++++++++++++++++++++++++++++++++++++++++++++

\begin{definition}
    Si \( G\) est un groupe et si \( g,h\in G\), nous notons \( [g,h]=ghg^{-1}h^{-1}\)\nomenclature[R]{\( [g,h]\)}{commutateur dans un groupe} le \defe{commutateur}{commutateur!dans un groupe} de \( g\) et \( h\). 
\end{definition}

L'élément neutre est toujours un commutateur : pour \( g=h \), \( [g,g]=ggg^{-1}g^{-1}=e \).

\begin{definition}      \label{DEFooBNLPooShKYXa}
    Le \defe{groupe dérivé}{dérivé!groupe}\index{groupe!dérivé} de \( G\) est le sous-groupe noté \( D(G)\)\nomenclature[R]{\( D(G)\)}{groupe dérivé} ou \( [G,G]\)\nomenclature[R]{\( [G,G]\)}{groupe dérivé} engendré\footnote{Définition \ref{DefooRDRXooEhVxxu}.} par les commutateurs.
\end{definition}
Autrement dit, \( D(G)\) est l'intersection de tous les sous-groupes de \( G\) contenant tous les commutateurs. Le groupe \( D(G)\) contient toujours au moins le neutre parce que c'est un groupe.

En vertu du lemme~\ref{LemFUIZooBZTCiy}, le groupe dérivé de \( G\) est l'ensemble des produits finis de commutateurs. C'est-à-dire que si \( S_m\) est l'ensemble des produits de \( m\) commutateurs, alors
\begin{equation}
    D(G)=\bigcup_{m=1}^{\infty}S_m.
\end{equation}

\begin{lemma}   \label{LemMMOCooDJJJhy}
    Le groupe dérivé est un sous-groupe caractéristique\footnote{Définition \ref{DEFooUXXTooCCLmQe}.}, et un sous-groupe normal\footnote{Définition \ref{DEFooNIIMooFkZgvX}.}.
\end{lemma}

\begin{proof}
    Il est évident que si \( \alpha\in\Aut(G)\) alors
    \begin{equation}
        \alpha\big( [g,h] \big)=\big[ \alpha(g),\alpha(h) \big],
    \end{equation}
    c'est-à-dire que \( D(G)\) est un sous-groupe caractéristique. En particulier si \( c\) est un commutateur, alors \( xcx^{-1}\) en est encore un, ce qui montre que \( D(G)\) est normal dans \( G\). Plus spécifiquement,
    \begin{subequations}
        \begin{align}
        x(ghg^{-1}h^{-1})x^{-1}&=(xgx^{-1})(xhx^{-1})(xg^{-1}x^{-1})(xh^{-1}x^{-1})\\
        &=(xgx^{-1})(xhx^{-1})(xgx^{-1})^{-1}(xhx^{-1})^{-1}.
        \end{align}
    \end{subequations}
\end{proof}

\begin{proposition}\label{PropAPRGooHBkELf}
    Le groupe quotient \( G/D(G)\) est abélien.
\end{proposition}

\begin{proof}
    En ce qui concerne le fait que \( G/D(G)\) soit abélien, nous savons que pour tout \( g,h\in G\) nous avons \( h^{-1}g^{-1}hg\in D(G)\) et donc
    \begin{equation}
        [g][h]=[gh]=[ghh^{-1}g^{-1}hg]=[hg]=[h][g].
    \end{equation}
\end{proof}

Le groupe quotient \( G/D(G)\) est appelé l'\defe{abélianisé}{abélianisé} de \( G\) et est parfois noté \( G^{ab}\)\nomenclature[R]{\( G^{ab}\)}{groupe abélianisé de \( G\)}.

Si \( f\colon G\to A\) est un homomorphisme entre le groupe \( G\) et un groupe abélien \( A\), alors \( f\big( D(G) \big)=\{ 0 \}\). Du coup \( f\) passe au quotient de \( G\) par \( D(G)\), et il existe une unique application \( \bar f\colon G/D(G)\to A\) telle que \( f=\bar f\circ \pi\) où \( \pi\colon G\to G/D(G)\) est la projection canonique.



%+++++++++++++++++++++++++++++++++++++++++++++++++++++++++++++++++++++++++++++++++++++++++++++++++++++++++++++++++++++++++++
\section{Suite de composition}
%+++++++++++++++++++++++++++++++++++++++++++++++++++++++++++++++++++++++++++++++++++++++++++++++++++++++++++++++++++++++++++
\index{sous-groupe!normal}\index{groupe!quotient}\index{quotient!de groupe}

%TODO : citer la page de la wikiversité sur Jordan-Hölder.
%TODO : donner la définition d'un raffinement de suite de composition.

\begin{definition}  \label{DefJWZSooNcntfK}
Une \defe{suite de composition}{composition!suite de}\index{suite de composition} pour un groupe \( G\) est une suite finie de sous-groupes \( (G_i)_{i=0,\ldots, n}\) telle que
\begin{equation}
    \{ e \}=G_n\subseteq G_{n-1}\subseteq\ldots\subseteq G_1\subseteq G_0=G
\end{equation}
et telle que \( G_{i+1}\) est normal\footnote{Nous rappelons au cas où que «normal» signifie «distingué».} dans \( G_i\). Les groupes \( G_i/G_{i+1}\) sont les \defe{quotients}{quotient!dans une suite de composition} de la suite de composition.

    Une suite de \defe{Jordan-Hölder}{suite!de Jordan-Hölder}\index{Jordan-Hölder} est une suite de composition dont tous les quotients sont simples.
\end{definition}
L'objet de nos prochaines pérégrinations mathématiques est de montrer que tout groupe fini admet une suite de Jordan-Hölder (théorème~\ref{ThoLgxWIC}).

\begin{lemma}[du papillon\cite{NjCCfW}]\label{LemsKpXCG}
    Soit \( G\) un groupe et des sous-groupes \( A\) et \( B\). Soit \( A'\) normal dans \( A\) et \( B'\) normal dans \( B\). Alors
    \begin{enumerate}
        \item
            \( A'(A\cap B')\) est normal dans \( A'(A\cap B)\)
        \item
            \( (A'\cap B)B'\) est normal dans \( (A\cap B)B'\)
        \item
            Nous avons les isomorphismes de groupes
            \begin{equation}
                \frac{ A'(A\cap B) }{ A'(A\cap B') }\simeq\frac{ (A\cap B)B' }{ (A'\cap B)B' }\simeq\frac{ B'(A\cap B) }{ B'(A'\cap B) }.
            \end{equation}
    \end{enumerate}
\end{lemma}

\begin{proof}
    Nous n'allons pas démontrer chacun des points; pour plus de détails, nous dirons simplement que «la preuve est très similaire dans les autres cas».

    Commençons par montrer que \( A'(A\cap B')\) est un groupe. Si \( a,b\in A'\) et \( x,y\in A\cap B'\),
    \begin{equation}
        axby=xx^{-1}axbx^{-1}xy
    \end{equation}
    En utilisant la normalité, \( x^{-1}ax\in A'\), donc \( xx^{-1}axbx^{-1}\in A'\) et donc le tout est dans \( A'(A\cap B')\). L'ensemble \( A'(A\cap B')\) est également stable pour l'inverse parce que
    \begin{equation}
        x^{-1}a^{-1}=\underbrace{x^{-1}a^{-1}x}_{\in A'}x^{-1}.
    \end{equation}

    Nous montrons maintenant que \( A'(A\cap B')\) est normal dans \( A'(A\cap B)\). Soient \( a,b\in A'\), \( x\in A\cap B'\) et \( f\in A\cap B\). Alors
    \begin{subequations}
        \begin{align}
        (bf)^{-1}(ax)(bf)&=(bf)^{-1}(a\underbrace{xbx^{-1}}_{=c\in A'}xf)\\
        &=f^{-1}b^{-1}acxf\\
        &=f^{-1}b^{-1}acf\underbrace{f^{-1}xf}_{=y\in A\cap B'}\\
        &=\underbrace{f^{-1}b^{-1}acf}_{\in A'}y\\
        &\in A'(A\cap B').
        \end{align}
    \end{subequations}

    Pour prouver l'isomorphisme
    \begin{equation}
        \frac{ A'(A\cap B) }{ A'(A\cap B') }=\frac{ (A\cap B)B' }{ (A'\cap B)B' },
    \end{equation}
    nous allons utiliser le deuxième théorème d'isomorphisme (\ref{ThoezgBep}\ref{ItembgDQEN}). Que nous appliquons à \( H=A\cap B\) et \( N=A'(A\cap B')\). La vérification que \( H\) normalise \( N\) est usuelle. Nous commençons par écrire
    \begin{equation}    \label{EqkphNsE}
        \frac{ A'(A\cap B')(A\cap B) }{ A'(A\cap B') }\simeq\frac{ A\cap B }{ A\cap B\cap A'(A\cap B') }.
    \end{equation}
    Pour simplifier un peu cette expression nous prouvons d'abord que
    \begin{equation}    \label{EqkhsyNh}
        (A\cap B)\cap A'(A\cap B')=(A'\cap B)(A\cap B').
    \end{equation}
    L'inclusion \( \supset\) est facile. Pour l'autre sens, étant donné que \( A'(A\cap B')\subset A\) nous avons
    \begin{equation}
        A\cap B\cap A'(A\cap B)=B\cap A'(A\cap B).
    \end{equation}
    Un élément de \( B\cap A'(A\cap B)\) est un élément de \(   B\) qui s'écrit sous la forme \( s=ax\) avec \( a\in A'\) et \( x\in A\cap B'\). Nous avons alors \( a=sx^{-1}\) avec \( s\in B\) et \( x^{-1} \in B'\). Par conséquent \( a\in B\) et donc \( a\in A'\cap B\). Nous avons donc
    \begin{equation}
        (A\cap B)\cap A'(A\cap B')=B\cap A'(A\cap B)\subset (A'\cap B)(A\cap B'),
    \end{equation}
    et donc l'égalité \eqref{EqkhsyNh}. Toujours dans l'idée de simplifier \eqref{EqkphNsE} nous remarquons que \( A\cap B'\) est un sous-ensemble de \( A\cap B'\), donc \( A'(A\cap B')(A\cap B)=A'(A\cap B)\). Il reste donc
    \begin{equation}
        \frac{ A'(A\cap B) }{ A'(A\cap B') }=\frac{ A\cap B }{ (A'\cap B)(A\cap B') }.
    \end{equation}
    Étant donné que les hypothèses sur \( A\) et \( B\) sont symétriques, le membre de droite peut aussi s'écrire en inversant \( A\) et \( B\). Nous en sommes à
    \begin{equation}
        \frac{ B'(A\cap B) }{ B'(A'\cap B) }=\frac{ A'(A\cap B) }{ A'(A\cap B') }.
    \end{equation}
    Nous devons encore justifier \( B'(A\cap B)=(A\cap B)B'\) et \( B'(A'\cap B)=(A'\cap B)B'\). Faisons le premier et laissons le second \href{http://abstrusegoose.com/395}{au lecteur}.
    Si \( b\in B'\) et \( x\in A\cap B\), alors
    \begin{equation}
        bx=x\underbrace{x^{-1}bx}_{\in B'}\in (A\cap B)B'.
    \end{equation}
\end{proof}

\begin{proposition}
    Si \( G\) est un groupe fini et que \( (G_i)\) est une suite de composition pour \( G\), alors l'ordre de \( G\) est le produit des ordres de ses quotients.
\end{proposition}

\begin{proof}
    Étant donné que \( G_{i+1}\) est toujours normal dans \( G_i\), le théorème de Lagrange (\ref{ThoLagrange}) s'applique et nous avons à chaque pas de la suite de composition nous avons
    \begin{equation}
        | \frac{ G_i }{ G_{i+1} } |=\frac{ | G_i | }{ | G_{i+1} | }
    \end{equation}
    et il suffit d'écrire \( | G |\) de façon télescopique :
    \begin{equation}
        | G |=\prod_{0\leq i\leq n-1}\frac{ | G_i | }{ | G_{i+1} | }
    \end{equation}
\end{proof}

Nous disons que les deux suites de composition \( (G_i)_{0\leq i\leq r}\) et \( (G_j)_{0\leq j\leq s}\) sont \defe{équivalentes}{equivalence@équivalence!suite de composition} si \( r=s\) et s'il existe une permutation \( \sigma\in S_{r-1}\) telle que
\begin{equation}
    \frac{ G_i }{ G_{i+1} }\simeq\frac{ H_{\sigma(i)} }{ H_{\sigma(i)+1} }.
\end{equation}

\begin{proposition}[Schreider]\index{lemme!de Schreider}
    Deux suites de composition d'un même groupe admettent des raffinements équivalents.
\end{proposition}

\begin{proof}
    Soient les suites de composition
    \begin{subequations}
        \begin{align}
            \{ e \}=G_m\subseteq\ldots\subseteq G_1\subseteq G_0=G\\
            \{ e \}=H_m\subseteq\ldots\subseteq H_1\subseteq H_0=G
        \end{align}
    \end{subequations}
    Nous raffinons la suite \( (G_i)\) en remplaçant \( G_{i+1}\subseteq G_i\) par
    \begin{equation}
        G_{i+1}=G_{i+1}(G_i\cap H_n)\subset G_{i+1}(G_i\cap H_{n-1})\subseteq\ldots\subseteq G_{i+1}(G_i\cap H_0)=G_i,
    \end{equation}
    et de même pour \( (H_j)\). Le groupe \( G_{i+1}(G_i\cap H_k)\) est normal dans \( G_{i+1}(G_i\cap H_{k-1})\) parce que \( G_{i+1}\) étant normal dans \( G_i\) et \( H_k\) dans \( H_{k-1}\), le lemme~\ref{LemsKpXCG} s'applique. Nous avons donc bien défini un raffinement.

    Nous devons maintenant prouver que les deux raffinements ainsi construits sont des suites de composition équivalentes. D'abord elles ont la même longueur \( mn\) parce que chacun des \( m\) éléments de la suite \( (G_i)\) a été remplacé par \( n\) éléments et inversement, chacun de \( n\) éléments de la suite \( (H_j)\) a été remplacé par \( m\) éléments.

    Par ailleurs, les quotients du raffinement de \( (G_i)\) sont de la forme
    \begin{equation}    \label{EqPAYTCB}
        \frac{ G_{i+1}(G_i \cap H_k) }{ G_{i+1}(G_i\cap H_{k+1}) }\simeq \frac{ H_{k+1}(H_k\cap G_i) }{ H_{k+1}(H_k\cap G_{i+1}) }
    \end{equation}
    en vertu du lemme du papillon (\ref{LemsKpXCG}). Le membre de droite de \eqref{EqPAYTCB} est un des quotients du raffinement de \( (H_j)\).
\end{proof}

\begin{lemma}[Schreider strictement décroissant]    \label{LemBSicRJ}
    Soient \( \Sigma_1\) et \( \Sigma_2\), deux suites de composition strictement décroissantes du groupe \( G\). Alors elles admettent des raffinements équivalents strictement décroissants.
\end{lemma}

\begin{proof}
    Par hypothèse, \( \Sigma_1\) et \( \Sigma_2\) n'ont pas de répétitions. Soient \( \Sigma''_1\) et \( \Sigma''_2\), des raffinements équivalents donnés par le lemme de Schreider. Étant donné que ce sont des suites de composition équivalentes, elles ont le même nombre de quotients réduits à \( \{ e \}\), c'est-à-dire le même nombre de répétitions.

    Les suites \( \Sigma'_1\) et \( \Sigma'_2\) obtenues en retirant les répétitions de \( \Sigma''_1\) et \( \Sigma''_2\) sont des raffinements équivalents de \( \Sigma_1\) et \( \Sigma_2\) et strictement décroissants.
\end{proof}

\begin{theorem}[Jordan-Hölder]\label{ThoLgxWIC}
    Tout groupe fini admet une suite de Jordan-Hölder.

    Deux suites de Jordan-Hölder sont équivalentes.
\end{theorem}
% TODO : trouver une preuve du fait que tout groupe fini admet une suite de Jordan-Hölder.

\begin{proof}
    Nous ne prouvons que le second point.

    Par définition, une suite de Jordan-Hölder n'a pas de raffinement strictement décroissant (à part elle-même) parce que \( G_{i+1}\) est normal maximum dans \( G_i\). Si \( \Sigma_1\) et \( \Sigma_2\) sont des suites de Jordan-Hölder nous pouvons considérer les raffinements équivalents strictement décroissants \( \Sigma'_1\) et \( \Sigma'_2\) du lemme de Schreider~\ref{LemBSicRJ}. Nous avons \( \Sigma'_1\sim\Sigma'_2\), mais par ce que nous venons de dire à propos de la maximalité, \( \Sigma'_1=\Sigma_1\) et \( \Sigma'_2=\Sigma_2\). D'où le résultat.
\end{proof}

%+++++++++++++++++++++++++++++++++++++++++++++++++++++++++++++++++++++++++++++++++++++++++++++++++++++++++++++++++++++++++++
\section{Groupes résolubles}
%+++++++++++++++++++++++++++++++++++++++++++++++++++++++++++++++++++++++++++++++++++++++++++++++++++++++++++++++++++++++++++

\begin{definition}  \label{DefOSYNooTROIKs}
    Le groupe \( G\) est \defe{résoluble}{groupe!résoluble} s'il existe une suite finie de sous-groupes \( G_i\)
    \begin{equation}
        \{ e \}=G_n\subset G_{n-1}\subset\ldots\subset G_1\subset G_0=G
    \end{equation}
    avec \( G_i\) normal dans \( G_{i+1}\) et \( G_i/G_{i+1}\) abélien.
\end{definition}
Il s'agit d'un groupe qui admet une suite de composition\footnote{Voir définition~\ref{DefJWZSooNcntfK}.} dont les quotients sont abéliens.

\begin{lemma}[\cite{HQRooKGAfpu}]   \label{LemOARMooYhYmbH}
    Soit \( G\) un groupe et \( H\) un sous-groupe normal. Le groupe \( G/H\) est abélien si et seulement si \( D(G)\subset H\).
\end{lemma}

\begin{proof}
    Les propositions suivantes sont équivalentes :
    \begin{itemize}
        \item Le groupe \( G/H\) est abélien
        \item pour tout \( x,y\in G\), \( \bar x\bar y=\bar y\bar x\)
        \item $\bar x\bar y\bar x^{-1}\bar y^{-1}=\bar e$
        \item \( \overline{ xyx^{-1}y^{-1} }=\bar e\)
        \item \( [x,y]\in H\)
        \item \( D(G)\subset H\).
    \end{itemize}
\end{proof}

\begin{proposition}[\cite{HQRooKGAfpu}] \label{PropRWYZooTarnmm}
    Un groupe est résoluble si et seulement si sa suite dérivée termine sur \( \{ e \}\).
\end{proposition}

\begin{proof}
    Grâce au lemme~\ref{LemMMOCooDJJJhy} et à la proposition~\ref{PropAPRGooHBkELf}, si la suite dérivée termine sur \( \{ e \}\) alors la suite dérivée est une suite qui répond aux conditions de la définition~\ref{DefOSYNooTROIKs} de groupe résoluble.

    Il faut donc encore montrer le sens direct. Nous supposons que \( G\) est un groupe résoluble et nous étudions sa suite dérivée. Nous avons une suite
    \begin{equation}
        \{ e \}=G_n\subset G_{n-1}\subset\ldots\subset G_1\subset G_0=G
    \end{equation}
    avec \( G_i/G_{i+1}\) abélien et \( G_{i+1}\) normal dans \( G_i\). Nous allons prouver par récurrence que \( D^i(G)\subset G_i\).

    Pour \( i=0\) nous avons bien \( G\subset G_0\). Notre hypothèse de récurrence est :
    \begin{equation}    \label{EqEAQEooEaeIEo}
        D^i(G)\subset G_i
    \end{equation}
    Par le lemme~\ref{LemOARMooYhYmbH} nous avons aussi
    \begin{equation}    \label{EqEDJXooLOLQcr}
        D(G_i)\subset G_{i+1}.
    \end{equation}
    En dérivant \eqref{EqEAQEooEaeIEo} et en tenant compte de \eqref{EqEDJXooLOLQcr}, \( D^{i+1}(G)\subset D(G_i)\subset G_{i+1}\). Donc par récurrence nous avons bien \( D^k(G)\subset G_k\) pour tout \( k\). Mais \( G_r=\{ e \}\) pour un certain \( r\), donc pour ce \( r\) nous avons \( D^r(G)=\{ e \}\), ce qu'il fallait.
\end{proof}

\begin{proposition} \label{PropBNEZooJMDFIB}
    Soient des groupes \( G\) et \( H\). Nous supposons que \( G\) est résoluble et nous considérons un homomorphisme \( \psi\colon G\to H\). Alors \( \psi(G)\) est résoluble.
\end{proposition}

\begin{proof}
    Vu que \( G\) est résoluble, il existe une suite de sous-groupes \( G_i\) tels que
    \begin{equation}
        \{ e \}=G_n\subset G_{n-1}\subset\ldots\subset G_1\subset G_0=G
    \end{equation}
    avec \( G_i\) normal dans \( G_{i+1}\) et \( G_i/G_{i+1}\) abélien. Nous posons \( \psi(G)_i=\psi(G_i)\) et nous avons \( \psi(G)_n=\psi\big( \{ e \} \big)=\{ e \}\) ainsi que \( \psi(G)_0=\psi(G)\); donc
    \begin{equation}
        \{ e \}=\psi(G)_n\subset \psi(G)_{n-1}\subset\ldots\subset \psi(G)_1\subset \psi(G)_0=\psi(G).
    \end{equation}

    Les faits que \( \psi(G)_i\) soit normal dans \( \psi(G)_{i+1}\) et que \( \psi(G)_i/\psi(G)_{i+1}\) soit abélien est directement la proposition~\ref{PropSRMJooIDPBoW}.

\end{proof}

%+++++++++++++++++++++++++++++++++++++++++++++++++++++++++++++++++++++++++++++++++++++++++++++++++++++++++++++++++++++++++++
\section{Action de groupes}
%+++++++++++++++++++++++++++++++++++++++++++++++++++++++++++++++++++++++++++++++++++++++++++++++++++++++++++++++++++++++++++

\begin{definition}[Thème~\ref{THEMEooKZHBooRCULcr}]  \label{DefActionGroupe}
    Une \defe{action de groupe}{action}\index{action} \( G\) sur un ensemble \( E\) est la donnée, pour chaque élément \( g \in G\), d'une fonction \(\phi_g : E \to E \), de telle sorte que:
    \begin{gather*}
        \phi_{e}(x) = x, \hspace{2em} \forall x \in E;\\
        \phi_{gh}(x) = \phi_g (\phi_h (x)),  \hspace{2em} \forall g,h \in G, \forall x \in E.
     \end{gather*}
     On dit dans ce cas que \( G \) \defe{agit}{action} sur \( E \).
\end{definition}

\begin{lemma}
    Pour tout \( g\in G\),
    \begin{enumerate}
        \item
            L'application \( \phi_g\colon E\to E\) est injective,
        \item
            Pour l'inverse : \( (\phi_g)^{-1}=\phi_{g^{-1}}\).
    \end{enumerate}
\end{lemma}

\begin{proof}
    Si \( x,y\in E\) sont tels que \( \phi_g(x)=\phi_g(y)\) alors en appliquant \( \phi_{g^{-1}}\) aux deux membres nous trouvons
    \begin{equation}
        (\phi_{g^{-1}}\phi_g)(x)=(\phi_{g^{-1}}\phi_g)(y),
    \end{equation}
    ce qui donne \( x=y\) parce que \( \phi_{g^{-1}}\phi_g=\phi_{g^{-1}g}=\phi_e=\id\).

    Les dernières trois égalités écrites disent que \( \phi_{g^{-1}}\) est l'inverse\footnote{Si vous décidez de dire ça a un jury dans un concours, soyez prêts à préciser les domaines.} de \( \phi_g\).
\end{proof}

Pour alléger les notations, on convient d'écrire $g \cdot x$, voire plus simplement $gx$ au lieu de \( \phi_g(x) \). Le deuxième axiome d'action de groupe dit que la notation $ghx$ ne souffre d'aucune ambiguïté.

\begin{definition}[Quelques notions autour de l'action]

    Si \( G\) agit sur un ensemble \( E\), nous notons \( G\cdot x\) l'\defe{orbite}{orbite!d'un point sous une action} de \( x\in E\) sous l'action de $G$:
\begin{equation*}
    G\cdot x = \{ gx \tq g \in G\}.
\end{equation*}

Nous notons \( G_x\) ou \( \Fix(x)\) le \defe{stabilisateur}{stabilisateur} de \( x\) :
\begin{equation}
    \Fix(x)=G_x=\{ g\in G\tq g\cdot x=x \}.
\end{equation}

    Pour \( g\in G\), nous notons enfin \( \Fix(g)\) le \defe{fixateur}{fixateur} de \( g\) :
\begin{equation}
    \Fix(g)=\{ x\in E\tq g\cdot x=x \}.
\end{equation}

\end{definition}

\begin{definition}  \label{DefuyYJRh}
    L'action de \( G\) sur \( E\) est \defe{fidèle}{fidèle (action)}\index{action!fidèle} si l'identité est le seul élément de \( G\) à fixer tous les points de \( E\), c'est-à-dire si \( gx=x\,\forall x\in E\Rightarrow g=e\).
\end{definition}

Un exemple d'action fidèle tout à fait non trivial sera donné avec l'action du groupe modulaire sur le plan de Poincaré dans le théorème~\ref{ThoItqXCm}.

Le groupe \( G\) agit toujours sur lui même à gauche et à droite. L'action à gauche est \( g\cdot h=gh\); celle à droite est \( g\cdot h=hg^{-1}\).

\begin{definition}      \label{DEFooCORTooEeOLPT}
    L'action \defe{adjointe}{action!adjointe} définie par \( g\cdot h=ghg^{-1}\) est une manière pour un groupe d'agir sur lui-même par automorphismes. Cela est souvent noté \( \AD(g)h=ghg^{-1}\).
\end{definition}
En effet pour tout \( g\in G\), l'application \( \AD(g)\colon G\to G\) est un automorphisme de \( G\).

Si \( H\) est un sous-groupe de  \( G\), nous notons \( G/H\) le quotient de $G$ par la relation \( g\sim gh\) pour tout \( h\in H\). Lorsque la distinction est importante, nous noterons \( (G/H)_g\)\nomenclature[R]{$(G/H)_g$}{classes à gauche} pour les classes à gauche et \( (G/H)_d\) pour les classes à droite.

Nous avons une relation d'équivalence à gauche et une à droite. D'abord
\begin{equation}
    x\sim_g y\Leftrightarrow xh=y
\end{equation}
pour un certain \( h\in H\). Ensuite
\begin{equation}
    x\sim_d y\Leftrightarrow hx=y
\end{equation}
pour un certain \( h\in H\).

Le lemme suivant est une généralisation du théorème de Lagrange~\ref{ThoLagrange}.

\begin{lemma}
    L'ensemble \( (G/H)_g\) est fini si et seulement si l'ensemble \( (G/H)_d\) est fini. S'il en est ainsi, alors \( (G/H)_g\) et \( (G/H)_d\) ont même cardinal qui vaut l'indice de \( H\) dans \( G\).
\end{lemma}

\begin{proof}
    L'application
    \begin{equation}
        \begin{aligned}
            f\colon (G/H)_g&\to (G/H)_d \\
            [x]_g&\mapsto [x^{-1}]_d
        \end{aligned}
    \end{equation}
    est une bijection bien définie. En effet si \( x\sim_g y\), nous avons \( h\in H\) tel que \( y^{-1}h=x^{-1}\), c'est-à-dire que \( x^{-1}\sim_d y^{-1}\) et \( f\) est bien définie. Le fait que \( f\) soit surjective est évident. Pour l'injectivité, soient \( x, y \in G \) tels que
    \begin{equation}
        f([x]_g)=f([y]_g).
    \end{equation}
    Alors \( x^{-1}\sim_d y^{-1}\), ce qui implique l'existence de \( h\in H\) tel que \( hx^{-1}=y^{-1}\), ou encore que \( xh^{-1}=y\), ce qui signifie que \( x\sim_gy\).

    Pour l'énoncé à propos de l'indice, nous procédons en plusieurs étapes simples.
    \begin{enumerate}
        \item
            Les classes (les éléments de \( (G/H)_g\)) forment une partition de $G$.
        \item
            Toutes les classes ont le même nombre d'éléments par la bijection
            \begin{equation}
                \begin{aligned}
                    f\colon [x]_g&\to [y]_g \\
                    xh&\mapsto yh.
                \end{aligned}
            \end{equation}
        \item
            Le nombre d'éléments dans une classe est égal à \( | H |\) par la bijection
            \begin{equation}
                \begin{aligned}
                    g\colon [x]_g&\to H \\
                    xh&\mapsto h.
                \end{aligned}
            \end{equation}
    \end{enumerate}
    Par conséquent
    \begin{equation}
        | G |=| H |\cdot \text{nombre de classes}=| H |\cdot\text{cardinal de $(G/H)_g$},
    \end{equation}
    et nous avons bien
    \begin{equation}
        \text{cardinal de }(G/H)_g=\frac{ | G | }{ | H | }=| G:H |.
    \end{equation}
\end{proof}

\begin{proposition}[Orbite-stabilisateur\cite{Combes}]     \label{Propszymlr}
    Soit \( G\) un groupe agissant sur un ensemble \( E\) et \( x\in E\).
    \begin{enumerate}
        \item
            Les ensembles \( G\cdot x\) et \( G/G_x\) sont équipotents.
        \item       \label{ITEMooCWUGooCOFHYk}
            L'orbite de \(\Fix(x)\) est finie si et seulement si \( \Fix(x)\) est d'indice fini dans \( G\). Dans ce cas nous avons
            \begin{equation}        \label{EqnLCHCE}
                \Card(G\cdot x)=| G:\Fix(x) |.
            \end{equation}
            Une autre façon d'écrire la même formule :
            \begin{equation}        \label{EqCewSXT}
                | G |=| \Fix(x) | |\mO_x |.
            \end{equation}
    \end{enumerate}
\end{proposition}
\index{équation!orbite-stabilisateur}
C'est la formule \eqref{EqnLCHCE} qui est nommée \wikipedia{fr}{Action_de_groupe_(mathématiques)\#Formule_des_classes.2C_formule_de_Burnside}{formule des classes} sur wikipédia.

\begin{proof}
    \begin{enumerate}
        \item
    Soit l'application
    \begin{equation}
        \begin{aligned}
            \psi\colon G\cdot x&\to G/G_x \\
            a\cdot x&\mapsto [a].
        \end{aligned}
    \end{equation}
    Cette application est bien définie parce que si \( a\cdot x=b\cdot x\), alors il existe \( h\in G_x\) tel que \( b=ah\), et par conséquent \( [a]=[b]\). Cette application est une bijection et par conséquent \( G\cdot x\) est équipotent à \( G/G_x\).
    \item
        Soit \( y\in \mO_x\) et \( A_y=\{ g\in G\tq g\cdot x=y \}\). L'ensemble \( A_y\) est une classe à gauche de \( \Fix(x)\), par conséquent \( | A_y |=|\Fix(x)|\) pour tout \( y\in\mO_x\). Les \( A_y\) pour différents \( y\) sont disjoints et nous avons de plus
        \begin{equation}
            \bigcup_{y\in\mO_x}A_y=G.
        \end{equation}
        Les ensembles \( A_y\) divisent donc \( G\) en \( | \mO_x |\) paquets de \( | \Fix(x) |\) éléments. D'où la formule \eqref{EqCewSXT}.

    \end{enumerate}
\end{proof}

\begin{corollary}       \label{CORooRRVHooTyCjZZ}
    Soit \( C_g\) la classe de conjugaison d'un élément  \( g\) du groupe fini \( G\). Alors
    \begin{equation}
        \Card(C_g)=| G:Z_G(g) |
    \end{equation}
    où $Z_G(g)$ est le centralisateur de \( g\) dans \( G\)\footnote{Définition~\ref{defGroupeCentre}.} de \( G\).
\end{corollary}

\begin{proof}
    Cela est une application de la proposition~\ref{Propszymlr} (formule \eqref{EqnLCHCE}) dans le cas de l'action adjointe de \( G\) sur lui-même.

    En effet, si nous considérons l'action adjointe, l'orbite est la classe de conjugaison : \( C_g=G\cdot g\). Et le stabilisateur de \( g\) pour l'action adjointe n'est autre que le centralisateur de \( g\) :
    \begin{subequations}
        \begin{align}
        \Fix(g)&=\{ h\in G\tq h\cdot g=g \}\\
        &= \{ h\in G\tq hgh^{-1}=g \}\\
        &=\{ h\in G\tq gh=hg \}\\
        &=Z_G(g).
        \end{align}
    \end{subequations}

    Donc la formule \( \Card(G\cdot g)=| G:G_g |\) devient, dans le cas de l'action adjointe de \( G\) sur lui-même : \( \Card(C_g)=| G:Z_G(g) |\).
\end{proof}

\begin{lemma}
    Soit \( G\) un groupe agissant sur l'ensemble \( E\). On définit \( x\sim x'\) si et seulement s'il existe \( g\in G\) tel que \( g\cdot x=x'\). Alors
    \begin{enumerate}
        \item
            la relation \( \sim\) est une relation d'équivalence.
        \item
            la classe \( [x]\) est l'orbite \( \mO_x\) de \( x\) sous \( G\).
    \end{enumerate}
\end{lemma}

\begin{corollary}[Équation des orbites]\index{équation!des orbites} \label{CorARFVMP}
    Soit \( G\) un groupe agissant sur l'ensemble \( E\) et \( \mO_1,\ldots, \mO_k  \) la liste des orbites (distinctes). Alors
    \begin{enumerate}
        \item
            \( E=\bigcup_i\mO_i\), l'union est disjointe,
        \item
            \( \Card(E)=\sum_i\Card(\mO_i)\).
    \end{enumerate}
\end{corollary}

\begin{definition}  \label{DefcSuYxz}
    Soit \( G\) un groupe agissant sur l'ensemble \( E\). Un \defe{domaine fondamental}{domaine!fondamental d'une action}\index{fondamental!domaine d'une action}\index{action!domaine fondamental} ou une \defe{transversale}{transversale} est une partie de \( E\) contenant un et un seul élément de chaque orbite.
\end{definition}
Autrement dit, les images des éléments d'un domaine fondamental forment une partition de l'ensemble :
\begin{equation}
    E=\bigsqcup_{g\in G}g(F)
\end{equation}
où  \( g(F) = \phi_g(F) = \{ \phi_g(x) \tq x \in F\} \). L'union est disjointe, c'est-à-dire que si \( g\neq g'\), alors \( g(F)\cap g'(F)=\emptyset\).


%TODO:
% - définir ce qu'est un p-groupe, au bon endroit (à côté du théorème de Lagrange?).
% - pour l'instant, la définition de p-groupe est là où on l'utilise, c'est-à-dire un peu avant les théorèmes de Sylow.
% - Il vaut peut-être mieux créer un thème ``indice dans un groupe'' dans lequel on parlera entre autres du théorème de Lagrange et de p-groupe.


\begin{proposition}[Équation des classes\cite{FabricegPSFinis}]     \label{PropUyLPdp}
    Soit \( G\), un groupe fini opérant sur un ensemble \( E\). Si \( E''\) est un ensemble contenant exactement un élément de chaque orbite dans \( E\setminus\Fix_G(E)\), alors
    \begin{equation}        \label{EqobuzfK}
        | G |=| \Fix_G(E) |+\sum_{x\in E''}\frac{ | G | }{ | \Fix_G(x) | }.
    \end{equation}
    Si de plus \( G\) est un $p$-groupe, alors
    \begin{equation}    \label{EqbzLEVJ}
        | E |=| \Fix_G(E) |\mod p.
    \end{equation}
\end{proposition}

\begin{proof}
    Par le corolaire~\ref{CorARFVMP}, nous avons \( | G |=\sum_{x\in E'}| \mO_x |\) où \( E'\) est une transversale.  En séparant la somme entre les orbites à un élément et les autres,
    \begin{equation}    \label{EqeggkBs}
        | G |=\Card(\Fix_G(E))+\sum_{x\in E''}\frac{ | G | }{ | \Fix_G(x) | }
    \end{equation}  \label{EqDgYbhm}
    où nous avons utilisé le fait que \( | G |=| \Fix_G(x) | |\mO_x |\).

    Si \( G\) est un \( p\)-groupe alors si \( x\in E''\), \( \Fix_G(x)\) est un sous-groupe propre de \( G\) et donc \( | \Fix_G(x) |\) est un diviseur propre de \( | G |\). Du coup la fraction \( | G |/|\Fix_G(x)|\) est une puissance non nulle de \( p\) et l'équation \eqref{EqobuzfK} devient immédiatement \eqref{EqbzLEVJ}.
\end{proof}


\begin{corollary}[Équation des classes]\index{équation!des classes}
    Soit \( G\), un groupe et \( C_1\),\ldots, \( C_l\) la liste de ses classes de conjugaison contenant plus de un élément. Alors
    \begin{equation}        \label{EqkgGmoq}
        \Card(G)=\Card\big( Z(G) \big)+\sum_i| G:Z_{g_i} |=\Card\big( Z(G) \big)+\sum_i\frac{ \Card(G) }{ \Card\big( \Fix(g_i) \big) }
    \end{equation}
    si \( g_i\in C_i\).
\end{corollary}

\begin{proof}
    Étant donné que les classes de conjugaison sont disjointes, le cardinal de \( G\) est bien la somme des cardinaux de ses classes. Les classes ne contenant que un seul élément sont celles des éléments de \( Z(G)\). En ce qui concerne les autres orbites, \( \Card(C_{g_i})=| G:Z_{g_i} |\) par le théorème orbite-stabilisateur (proposition~\ref{Propszymlr}).
\end{proof}

\begin{theorem}[\wikipedia{fr}{Action_de_groupe_(mathématiques)}{Formule de Burnside}]      \label{THOooEFDMooDfosOw}
    Si \( G\) est un groupe fini agissant sur l'ensemble fini \( E\) et si \( \Omega\) est l'ensemble des orbites, alors
    \begin{equation}    \label{EqTUsblv}
        \Card(\Omega)=\frac{1}{ | G | }\sum_{g\in G}\Card\big( \Fix(g) \big).
    \end{equation}
\end{theorem}
\index{Burnisde!formule}
\index{formule!Burnside}

\begin{proof}
    Nous considérons l'ensemble
    \begin{equation}
        A=\{ (g,x)\in G\times E\tq gx=x \},
    \end{equation}
    et nous en calculons le cardinal de deux façons. D'abord
    \begin{subequations}
        \begin{align}
            \Card(A)&=\sum_{x\in E}\Card\{ g\in g\tq gx=x \}\\
            &=\sum_{x\in E}\Card(\Fix(x))\\
            &=\sum_{\omega\in \Omega}\sum_{x\in \omega}\Card(\Fix(x))\\
            &=\sum_{\omega\in \Omega}\frac{ | G | }{ \Card(\omega) }     \label{EqyVtkyf}\\
            &=| G |.
        \end{align}
    \end{subequations}
    Pour obtenir \eqref{EqyVtkyf} nous avons utilisé l'équation des classes \eqref{EqCewSXT}. L'autre façon de calculer \( \Card(A)\) est de regrouper ainsi :
    \begin{equation}
        \Card(A)=\sum_{g\in G}\Card\{ x\in E\tq gx=x \}=\sum_{g\in G}\Card(\Fix(g)).
    \end{equation}
    En égalisant les deux expressions de \( \Card(A)\) nous trouvons
    \begin{equation}
        | G |\Card(\Omega)=\sum_{g\in G}\Card(\Fix(g)).
    \end{equation}
\end{proof}

\begin{proposition}
    Soit \( G\) un groupe et \( H\), un sous-groupe du centre de \( G\).
    \begin{enumerate}
        \item
            \( H\) est normal dans \( G\).
        \item
            Si \( G/H\) est monogène, alors \( G\) est abélien.
        \item
            Si \( G\) est fini de centre \( Z\), alors \( | G:H |\) n'est pas premier.
    \end{enumerate}
\end{proposition}

% TODO: preuve

\begin{theorem}     \label{THOooRGSTooIWyhqt}
    Soit \( G\) un groupe cyclique\footnote{Définition \ref{DefHFJWooFxkzCF}.} d'ordre \( n\).
    \begin{enumerate}
        \item
            Tout sous-groupe de \( G\) est cyclique.
        \item
            Pour chaque diviseur \( d\) de \( n\), il existe un unique sous-groupe \( H_d\) de \( G\) d'ordre \( d\).
    \end{enumerate}
    Si \( a\) est un générateur de \( G\), alors \( H_d\) peut être décrit des façons suivantes :
    \begin{equation}
        H_d=\{ x\in G\tq x^d=e \}=\{ x\in G\tq\exists y\in G\tq y^{n/d}=x \}=\langle a^{n/d}\rangle.
    \end{equation}
\end{theorem}

% TODO: preuve

\begin{definition}      \label{DEFooQDHPooCfDEuL}
    Soit \( G\) un groupe agissant sur un ensemble \( E\). Nous disons que l'action est \defe{transitive}{transitive}\index{action!transitive} si elle possède une seule orbite. L'action est \defe{libre}{libre!action}\index{action!libre} si \( g\cdot x=g'\cdot x\) implique \( g=g'\).
\end{definition}
